\chapter{Accidentally Eating Unsafe Foods}

\section*{Definition and Overview}
Accidentally eating unsafe foods, commonly called food poisoning or foodborne illness, means becoming unwell after eating or drinking something contaminated. Contamination can be caused by bacteria, viruses, parasites, or by chemicals and natural toxins within the food. Illness most often affects the stomach and bowels, causing vomiting and diarrhoea, but some germs produce fever, bloody diarrhoea, or serious systemic illness, particularly in young children, older people, pregnant women, and those with weakened immune systems. Most foodborne illness is mild and settles within a few days, but a small proportion of cases are severe and require hospital care.

\section*{Epidemiology}
Foodborne illness is very common: most people experience it at some point in their lives, and in Australia it is estimated to affect several million people every year. Common sources include undercooked meat, poultry, and eggs, unpasteurised dairy products, raw or undercooked seafood, unwashed salads and sprouts, and leftovers that were not refrigerated properly. Norovirus and \textit{Campylobacter} are among the most frequent causes. Outbreaks commonly occur in settings where food is prepared in quantity, such as restaurants, childcare centres, and residential care facilities. The people most at risk of severe illness are infants, pregnant women, older adults, and those with weakened immune systems.

\section*{Aetiology and Pathophysiology}
Unsafe food makes people ill in two main ways: by infection and by poisoning. Infection occurs when living germs multiply in the gut or body; poisoning occurs when pre-formed toxins in food cause symptoms.

\begin{itemize}
    \item \textbf{Viruses:} norovirus and hepatitis A spread easily from person to person and through contaminated food and water
    \item \textbf{Bacteria:} \textit{Salmonella}, \textit{Campylobacter}, \textit{Escherichia coli}, \textit{Listeria}, \textit{Staphylococcus aureus}, and \textit{Clostridium} species
    \item \textbf{Parasites:} such as \textit{Giardia}, acquired from contaminated water or food
    \item \textbf{Toxins:} some bacteria release toxins in food before it is eaten, such as staphylococcal toxin in cream-filled foods, and some seafood carries natural toxins
\end{itemize}

The mechanism varies by organism. Some germs infect the lining of the bowel directly, causing inflammation, fever, and sometimes bloody diarrhoea. Others release toxins that trigger vomiting and watery diarrhoea quickly, often within hours of eating. A few bacteria, notably \textit{Listeria}, can spread beyond the bowel into the bloodstream, which is why they are dangerous in pregnancy and in immunocompromised people.

\section*{Clinical Features}
Symptoms depend on the organism and typically begin within hours to a few days of eating.

\begin{itemize}
    \item Nausea and vomiting, often starting suddenly
    \item Diarrhoea, which may be watery or bloody
    \item Abdominal cramps and pain
    \item Fever, headache, and aching muscles
    \item Signs of dehydration: dry mouth, dark or reduced urine, dizziness, and weakness
    \item Less commonly, double vision, drooping eyelids, or muscle weakness (possible botulism), or tingling and flushing with certain seafood toxins
\end{itemize}

\section*{Differential Diagnosis}
The symptoms of foodborne illness overlap with several other conditions.

\begin{itemize}
    \item Viral gastroenteritis acquired from other people rather than from food
    \item Appendicitis or other abdominal emergencies presenting with vomiting and pain
    \item Inflammatory bowel disease or other chronic bowel conditions
    \item Food allergy or intolerance, such as lactose intolerance
    \item Medication side effects, including from some antibiotics
    \item Poisoning from chemicals, alcohol, or medicines
\end{itemize}

\section*{When to Seek Medical Care}
A doctor should be consulted if symptoms are severe, persist beyond a couple of days, or occur in a person at higher risk, such as a young child, pregnant woman, older adult, or someone with a weakened immune system.

\textbf{Call 000 (or local emergency services) for:}
\begin{itemize}
    \item Signs of severe dehydration: no urine output, confusion, collapse, or an inability to keep down any fluid
    \item Blood in the vomit, or large amounts of blood in the stool
    \item High fever with shaking chills
    \item Difficulty breathing, double vision, or muscle weakness, which may indicate botulism
    \item Severe abdominal pain, or symptoms in an infant who appears floppy or unwell
\end{itemize}

\section*{Diagnosis and Investigations}
Most cases are managed without testing, but investigation is appropriate in severe, persistent, or outbreak situations.

\begin{itemize}
    \item \textbf{History:} what was eaten, when, and whether others who shared the meal are also unwell
    \item \textbf{Examination:} assessment of hydration status and the abdomen for signs of an acute surgical problem
    \item \textbf{Stool tests:} samples are tested for bacteria, viruses, or parasites in severe cases, when symptoms persist, or during an outbreak
    \item \textbf{Blood tests:} to check for dehydration, kidney dysfunction, or spread of infection
    \item \textbf{Public health reporting:} serious cases and suspected contaminated products are reported so that sources can be traced and removed
\end{itemize}

\section*{Management}
Management is predominantly supportive, with specific treatment reserved for particular organisms.

\begin{itemize}
    \item \textbf{Fluids:} the main treatment is replacing lost fluid, with oral rehydration drinks for moderate cases and intravenous fluids in hospital for severe dehydration
    \item \textbf{Solid foods:} return to bland, easily digested food as symptoms settle
    \item \textbf{Anti-nausea medicines:} may help persistent vomiting
    \item \textbf{Antibiotics:} needed for some bacterial infections such as severe \textit{Salmonella}, \textit{Listeria}, or \textit{Campylobacter}, but they are not used for most mild cases and can worsen some infections
    \item \textbf{Avoid:} medicines that slow the bowel are generally not recommended when there is bloody diarrhoea or fever
    \item \textbf{Hygiene:} careful hand washing to avoid spreading the illness to others
    \item \textbf{Special care:} close monitoring for children, pregnant women, and older adults
\end{itemize}

\section*{Complications}
\begin{itemize}
    \item Dehydration, which is most dangerous in young children and older adults
    \item Kidney damage, including haemolytic uraemic syndrome, a rare but serious complication of some \textit{Escherichia coli} infections
    \item Spread of bacteria beyond the bowel, causing bloodstream infection
    \item \textit{Listeria} infection in pregnancy, which can harm the unborn baby
    \item Reactive arthritis after some infections, especially \textit{Campylobacter} and \textit{Salmonella}
    \item Guillain--Barr\'e syndrome, a rare nerve disorder that can follow \textit{Campylobacter} infection
    \item Persistent bowel symptoms (post-infectious irritable bowel syndrome)
    \item Spread to household members or contamination of food prepared for others
\end{itemize}

\section*{Prognosis and Outcomes}
Most people recover fully within a few days. Norovirus usually settles in 1--3 days, \textit{Salmonella} within about a week, and \textit{Campylobacter} usually within a week or so. Recovery is generally complete even without treatment, provided fluids are maintained. Severe illness, complications, and death are uncommon in healthy adults and mostly affect the very young, the elderly, pregnant women, and immunocompromised people. When complications do occur, they can take weeks to resolve.

\section*{Prevention and Self-Care}
\begin{itemize}
    \item Wash hands with soap and water before preparing food and after handling raw meat, eggs, or the toilet
    \item Keep raw meat, poultry, and seafood separate from ready-to-eat foods
    \item Cook meat and poultry thoroughly and reheat leftovers until steaming hot
    \item Chill perishable foods promptly and do not leave them at room temperature for more than about two hours
    \item Wash fruit and vegetables before eating
    \item Avoid unpasteurised milk and dairy products, and avoid high-risk foods such as soft cheeses, deli meats, and raw sprouts during pregnancy
    \item Check use-by dates and discard food that smells or looks unsafe
    \item Stay home and avoid preparing food for others while you have vomiting or diarrhoea
\end{itemize}

\section*{Sources and Further Reading}
The following organisations provide reliable information on foodborne illness and food safety.

\begin{itemize}
    \item NHS. \textit{Food poisoning: symptoms and treatment.} https://www.nhs.uk/conditions/food-poisoning/
    \item Centers for Disease Control and Prevention. \textit{Food safety and foodborne illness prevention.} https://www.cdc.gov/foodsafety/
    \item World Health Organization. \textit{Food safety and foodborne diseases.} https://www.who.int/
    \item Mayo Clinic. \textit{Food poisoning: symptoms and causes.} https://www.mayoclinic.org/diseases-conditions/food-poisoning/
    \item MedlinePlus. \textit{Foodborne illness health information.} https://medlineplus.gov/foodborneillness.html
    \item Australian Government Department of Health. \textit{Food safety and public health.} https://www.health.gov.au/
\end{itemize}
